\textbf{Phần 1. Trắc nghiệm lựa chọn}
\setcounterpageref{ex}{0}

\begin{ex} %[5H8V47]
Một bể nước ngầm dạng hình hộp chữ nhật chứa được tối đa $8{,}5\text{ m}^3$ nước. Biết hiện tại trong bể đã có $4\,750\text{ dm}^3$ nước. Người ta cần bơm thêm bao nhiêu lít nước nữa để đầy bể?
\choice
{$3\,250\text{ lít}$}
{\True $3\,750\text{ lít}$}
{$3\,500\text{ lít}$}
{$4\,250\text{ lít}$}
\loigiai{
Đổi: $8{,}5\text{ m}^3 = 8\,500\text{ dm}^3 = 8\,500\text{ lít}$.\\
Số lít nước đã có là $4\,750\text{ dm}^3 = 4\,750\text{ lít}$.\\
Số lít nước cần bơm thêm là:
$$8\,500 - 4\,750 = 3\,750\text{ (lít)}.$$
\textbf{Đáp án của câu hỏi là B}
}
\end{ex}

\begin{ex} %[5H8V47]
Một khối gỗ có thể tích $1\text{ m}^3\text{ }250\text{ dm}^3$. Người ta cưa lấy $3$ thanh gỗ, mỗi thanh có thể tích $85\text{ dm}^3$. Thể tích phần gỗ còn lại là
\choice
{$995\text{ dm}^3$}
{$1\,005\text{ dm}^3$}
{\True $995\text{ dm}^3$}
{$1\,165\text{ dm}^3$}
\loigiai{
Đổi $1\text{ m}^3\text{ }250\text{ dm}^3 = 1\,250\text{ dm}^3$.\\
Thể tích của $3$ thanh gỗ đã cưa là:
$$85 \times 3 = 255\text{ (dm}^3\text{)}.$$
Thể tích phần gỗ còn lại là:
$$1\,250 - 255 = 995\text{ (dm}^3\text{)}.$$
\textbf{Đáp án của câu hỏi là C}
}
\end{ex}

\begin{ex} %[5H8V47]
Giá trị của biểu thức $3\text{ m}^3 - 1\,450\text{ dm}^3 + 0{,}8\text{ m}^3$ viết dưới dạng số thập phân với đơn vị mét khối là
\choice
{$2{,}45\text{ m}^3$}
{\True $2{,}35\text{ m}^3$}
{$2{,}25\text{ m}^3$}
{$3{,}35\text{ m}^3$}
\loigiai{
Đổi: $1\,450\text{ dm}^3 = 1{,}45\text{ m}^3$.\\
Biểu thức trở thành:
$$3\text{ m}^3 - 1{,}45\text{ m}^3 + 0{,}8\text{ m}^3 = 1{,}55\text{ m}^3 + 0{,}8\text{ m}^3 = 2{,}35\text{ m}^3.$$
\textbf{Đáp án của câu hỏi là B}
}
\end{ex}

\begin{ex} %[5H8V47]
Mỗi mét khối cát nặng $1{,}5$ tấn. Một xe tải chở được tối đa $6$ tấn hàng. Xe tải đó chở được nhiều nhất bao nhiêu đề-xi-mét khối cát để không bị quá tải?
\choice
{$400\text{ dm}^3$}
{$40\,000\text{ dm}^3$}
{\True $4\,000\text{ dm}^3$}
{$4\text{ dm}^3$}
\loigiai{
Thể tích cát tối đa xe chở được tính theo mét khối là:
$$6 : 1{,}5 = 4\text{ (m}^3\text{)}.$$
Đổi $4\text{ m}^3 = 4\,000\text{ dm}^3$.\\
\textbf{Đáp án của câu hỏi là C}
}
\end{ex}

\begin{ex} %[5H8V47]
Một bể bơi mini có thể tích $15\text{ m}^3$. Người ta dùng một máy bơm có công suất $1\,250\text{ lít/giờ}$. Hỏi sau $8$ giờ bơm liên tục thì lượng nước trong bể chiếm bao nhiêu phần trăm dung tích bể?
\choice
{$60\%$}
{\True $66{,}67\%$}
{$75\%$}
{$80\%$}
\loigiai{
Sau $8$ giờ, máy bơm được lượng nước là:
$$1\,250 \times 8 = 10\,000\text{ (lít)} = 10\text{ m}^3.$$
Tỉ số phần trăm lượng nước trong bể so với dung tích bể là:
$$10 : 15 \times 100\% \approx 66{,}67\%.$$
\textbf{Đáp án của câu hỏi là B}
}
\end{ex}

\begin{ex} %[5H8V47]
Người ta xếp các khối lập phương cạnh $1\text{ dm}$ thành một khối hộp chữ nhật có kích thước dài $2\text{ m}$, rộng $1{,}5\text{ m}$ và cao $1\text{ m}$. Hỏi cần dùng tất cả bao nhiêu khối lập phương nhỏ đó?
\choice
{$300$}
{\True $3\,000$}
{$30\,000$}
{$300\,000$}
\loigiai{
Thể tích của khối hộp chữ nhật là:
$$2 \times 1{,}5 \times 1 = 3\text{ (m}^3\text{)}.$$
Đổi $3\text{ m}^3 = 3\,000\text{ dm}^3$.\\
Mỗi khối lập phương nhỏ có thể tích là $1\text{ dm}^3$.\\
Số khối lập phương nhỏ cần dùng là:
$$3\,000 : 1 = 3\,000\text{ (khối)}.$$
\textbf{Đáp án của câu hỏi là B}
}
\end{ex}

\begin{ex} %[5H8V47]
Cho biết $0{,}075\text{ m}^3 = \ldots\ldots\text{ cm}^3$. Số thích hợp điền vào chỗ chấm là
\choice
{$750$}
{$7\,500$}
{\True $75\,000$}
{$750\,000$}
\loigiai{
Vì $1\text{ m}^3 = 1\,000\,000\text{ cm}^3$ nên:
$$0{,}075\text{ m}^3 = 0{,}075 \times 1\,000\,000\text{ cm}^3 = 75\,000\text{ cm}^3.$$
\textbf{Đáp án của câu hỏi là C}
}
\end{ex}

\begin{ex} %[5H8V47]
Một công trình cần $24\text{ m}^3$ bê tông tươi. Người ta dùng các xe trộn bê tông, mỗi xe chở được $3\,200\text{ dm}^3$ bê tông. Hỏi cần ít nhất bao nhiêu chuyến xe để chở hết lượng bê tông đó?
\choice
{$7\text{ chuyến}$}
{\True $8\text{ chuyến}$}
{$9\text{ chuyến}$}
{$10\text{ chuyến}$}
\loigiai{
Đổi: $24\text{ m}^3 = 24\,000\text{ dm}^3$.\\
Thực hiện phép chia:
$$24\,000 : 3\,200 = 7{,}5.$$
Vì số chuyến xe phải là số tự nhiên nên cần ít nhất $8$ chuyến xe để chở hết lượng bê tông.
\textbf{Đáp án của câu hỏi là B}
}
\end{ex}

\begin{ex} %[5H8V47]
Biết $\dfrac{3}{4}\text{ m}^3$ nước được chia đều vào $6$ thùng chứa giống nhau. Thể tích nước trong mỗi thùng là
\choice
{$120\text{ dm}^3$}
{\True $125\text{ dm}^3$}
{$150\text{ dm}^3$}
{$175\text{ dm}^3$}
\loigiai{
Đổi: $\dfrac{3}{4}\text{ m}^3 = 0{,}75\text{ m}^3 = 750\text{ dm}^3$.\\
Thể tích nước trong mỗi thùng là:
$$750 : 6 = 125\text{ (dm}^3\text{)}.$$
\textbf{Đáp án của câu hỏi là B}
}
\end{ex}

\begin{ex} %[5H8V47]
Trong các số đo sau, số đo nào lớn nhất?
\choice
{$4\text{ m}^3\text{ }50\text{ dm}^3$}
{$4\,500\text{ dm}^3$}
{$4\,050\,000\text{ cm}^3$}
{\True $4{,}52\text{ m}^3$}
\loigiai{
Đổi các số đo về cùng đơn vị mét khối:
- $4\text{ m}^3\text{ }50\text{ dm}^3 = 4{,}05\text{ m}^3$.\\
- $4\,500\text{ dm}^3 = 4{,}5\text{ m}^3$.\\
- $4\,050\,000\text{ cm}^3 = 4{,}05\text{ m}^3$.\\
- $4{,}52\text{ m}^3$.\\
Vì $4{,}52 > 4{,}5 > 4{,}05$ nên số đo lớn nhất là $4{,}52\text{ m}^3$.\\
\textbf{Đáp án của câu hỏi là D}
}
\end{ex}

\begin{ex} %[5H8V47]
Một bể nước hình hộp chữ nhật có diện tích đáy là $4{,}5\text{ m}^2$. Khi mở vòi chảy vào bể được $5\,400\text{ lít}$ nước thì mực nước trong bể dâng lên thêm bao nhiêu mét?
\choice
{$0{,}8\text{ m}$}
{\True $1{,}2\text{ m}$}
{$1{,}5\text{ m}$}
{$0{,}12\text{ m}$}
\loigiai{
Đổi: $5\,400\text{ lít} = 5\,400\text{ dm}^3 = 5{,}4\text{ m}^3$.\\
Chiều cao mực nước dâng lên bằng thể tích nước chia cho diện tích đáy:
$$5{,}4 : 4{,}5 = 1{,}2\text{ (m)}.$$
\textbf{Đáp án của câu hỏi là B}
}
\end{ex}

\begin{ex} %[5H8V47]
Trung bình mỗi ngày một hộ gia đình dùng hết $450\text{ lít}$ nước. Hỏi trong tháng $4$ (có $30$ ngày), hộ gia đình đó đã sử dụng bao nhiêu mét khối nước?
\choice
{$135\text{ m}^3$}
{\True $13{,}5\text{ m}^3$}
{$1{,}35\text{ m}^3$}
{$1\,350\text{ m}^3$}
\loigiai{
Tổng số lít nước hộ gia đình đó dùng trong tháng $4$ là:
$$450 \times 30 = 13\,500\text{ (lít)}.$$
Đổi: $13\,500\text{ lít} = 13\,500\text{ dm}^3 = 13{,}5\text{ m}^3$.\\
\textbf{Đáp án của câu hỏi là B}
}
\end{ex}

\textbf{Phần 2. Trắc nghiệm đúng sai}
\setcounterpageref{ex}{0}

\begin{ex} %[5H8V47]
Cho hai số đo thể tích: $V_1 = 3{,}04\text{ m}^3$ và $V_2 = 3\text{ m}^3\text{ }40\text{ dm}^3$. Xét tính đúng/sai của các nhận định sau:
\choiceTF
{\True $V_1$ đọc là ba phẩy không bốn mét khối}
{\True $V_2$ đổi ra mét khối bằng $3{,}04\text{ m}^3$}
{Hiệu thể tích giữa $V_1$ và $V_2$ bằng $360\text{ dm}^3$}
{\True Tổng thể tích của $V_1$ và $V_2$ bằng $6\,080\text{ lít}$}
\loigiai{
\begin{itemchoice}
\item Đúng. $3{,}04\text{ m}^3$ đọc là ba phẩy không bốn mét khối.
\item Đúng. Vì $40\text{ dm}^3 = 0{,}04\text{ m}^3$ nên $V_2 = 3 + 0{,}04 = 3{,}04\text{ m}^3$.
\item Sai. Vì $V_1 = V_2 = 3{,}04\text{ m}^3$ nên hiệu $V_1 - V_2 = 0$.
\item Đúng. Tổng hai thể tích là $3{,}04 + 3{,}04 = 6{,}08\text{ m}^3 = 6\,080\text{ dm}^3 = 6\,080\text{ lít}$.
\end{itemchoice}
}
\end{ex}

\begin{ex} %[5H8V47]
Một bể chứa nước dạng hình hộp chữ nhật có dung tích $12\text{ m}^3$. Người ta bơm nước vào bể bằng một vòi có lưu lượng $800\text{ dm}^3$ mỗi giờ. Xét tính đúng/sai của các nhận định sau:
{\choiceTF
{\True Dung tích của bể bằng $12\,000\text{ lít}$}
{\True Sau $5$ giờ bơm, vòi bơm được $4\text{ m}^3$ nước}
{Sau $10$ giờ bơm thì nước tràn ra khỏi bể}
{\True Cần tổng cộng $15$ giờ bơm để nước đầy bể}
}
\loigiai{
\begin{itemchoice}
\item Đúng. $12\text{ m}^3 = 12\,000\text{ dm}^3 = 12\,000\text{ lít}$.
\item Đúng. Sau $5$ giờ, lượng nước bơm được là $800 \times 5 = 4\,000\text{ dm}^3 = 4\text{ m}^3$.
\item Sai. Sau $10$ giờ, lượng nước bơm được là $800 \times 10 = 8\,000\text{ dm}^3 = 8\text{ m}^3 < 12\text{ m}^3$ nên nước chưa tràn khỏi bể.
\item Đúng. Thời gian để bơm đầy bể là $12\,000 : 800 = 15\text{ (giờ)}$.
\end{itemchoice}
}
\end{ex}

\textbf{Phần 3. Câu trả lời ngắn}
\setcounterpageref{ex}{0}

\begin{ex} %[5H8V47]
Một bể nước có thể tích $5\text{ m}^3$. Hiện lượng nước trong bể đang chiếm $68\%$ thể tích bể. Hỏi cần đổ thêm vào bể bao nhiêu lít nước nữa để đầy bể?
\shortans{1600}
\loigiai{
Tỉ số phần trăm lượng nước cần đổ thêm là:
$$100\% - 68\% = 32\%.$$
Đổi $5\text{ m}^3 = 5\,000\text{ lít}$.\\
Số lít nước cần đổ thêm là:
$$5\,000 \times 32\% = 1\,600\text{ (lít)}.$$
}
\end{ex}

\begin{ex} %[5H8V47]
Người ta xây một bức tường gạch có thể tích $7{,}2\text{ m}^3$. Biết mỗi viên gạch chiếm thể tích $1\,200\text{ cm}^3$, và phần vữa trát vữa chiếm $20\%$ thể tích bức tường. Hỏi cần bao nhiêu viên gạch để xây bức tường đó?
\shortans{4800}
\loigiai{
Đổi: $7{,}2\text{ m}^3 = 7\,200\,000\text{ cm}^3$.\\
Thể tích vữa chiếm $20\%$ nên thể tích các viên gạch chiếm $100\% - 20\% = 80\%$ thể tích bức tường.\\
Tổng thể tích gạch là:
$$7\,200\,000 \times 80\% = 5\,760\,000\text{ (cm}^3\text{)}.$$
Số viên gạch cần dùng là:
$$5\,760\,000 : 1\,200 = 4\,800\text{ (viên)}.$$
}
\end{ex}

\begin{ex} %[5H8V47]
Một khối gỗ dạng hình lập phương cạnh $1\text{ m}$. Người ta cắt khối gỗ đó thành các khối lập phương nhỏ cạnh $2\text{ dm}$. Hỏi cắt được tất cả bao nhiêu khối lập phương nhỏ như vậy?
\shortans{125}
\loigiai{
Đổi $1\text{ m} = 10\text{ dm}$.\\
Số khối lập phương nhỏ theo mỗi cạnh là:
$$10 : 2 = 5\text{ (khối)}.$$
Tổng số khối lập phương nhỏ cắt được là:
$$5 \times 5 \times 5 = 125\text{ (khối)}.$$
(Hoặc: Thể tích khối gỗ lớn là $1\text{ m}^3 = 1\,000\text{ dm}^3$. Thể tích mỗi khối gỗ nhỏ là $2 \times 2 \times 2 = 8\text{ dm}^3$. Số khối là $1\,000 : 8 = 125$ khối).
}
\end{ex}

\begin{ex} %[5H8V47]
Một xe téc chở xăng chứa đầy $9\text{ m}^3$ xăng. Người ta đã bơm xăng từ xe téc đó vào $15$ thùng phuy giống nhau, mỗi thùng chứa $240\text{ dm}^3$. Hỏi trong xe téc còn lại bao nhiêu lít xăng?
\shortans{5400}
\loigiai{
Đổi $9\text{ m}^3 = 9\,000\text{ dm}^3 = 9\,000\text{ lít}$.\\
Lượng xăng đã bơm vào $15$ thùng phuy là:
$$240 \times 15 = 3\,600\text{ (dm}^3\text{)} = 3\,600\text{ (lít)}.$$
Số lít xăng còn lại trong xe téc là:
$$9\,000 - 3\,600 = 5\,400\text{ (lít)}.$$
}
\end{ex}

\begin{ex} %[5H8V47]
Một hố móng hình hộp chữ nhật có chiều dài $6\text{ m}$, chiều rộng $4\text{ m}$ và chiều sâu $2\text{ m}$. Người ta dùng xe tải chở đất đi đổ, mỗi chuyến xe chở được $3\text{ m}^3$ đất. Hỏi cần bao nhiêu chuyến xe để chở hết lượng đất đào từ hố móng đó?
\shortans{16}
\loigiai{
Thể tích đất đào từ hố móng là:
$$6 \times 4 \times 2 = 48\text{ (m}^3\text{)}.$$
Số chuyến xe cần để chở hết đất là:
$$48 : 3 = 16\text{ (chuyến)}.$$
}
\end{ex}

\begin{ex} %[5H8V47]
Một bể nước hình hộp chữ nhật đang chứa $1\,800\text{ dm}^3$ nước, chiếm $40\%$ thể tích của bể. Hỏi thể tích của bể nước đó là bao nhiêu đề-xi-mét khối?
\shortans{4500}
\loigiai{
Thể tích của bể nước là:
$$1\,800 : 40\% = 1\,800 \times 100 : 40 = 4\,500\text{ (dm}^3\text{)}.$$
}
\end{ex}
